% Síntesis axiomática incorporada desde el anexo de caminos.
% Las pruebas operatorias extensas permanecen en sus fuentes.

\section{Reconstruction of quantum principles from the genealogical state}
\label{sec:postulados-cuanticos-genealogicos-rev2}
\index{quantum mechanics!genealogical reconstruction}
\index{state!Hilbert realization}

Hilbert's formulation organizes quantum mechanics through states,
self-adjoint observables, unitary evolution, tensorial composition, and a
probabilistic rule for outcomes~\cite{vonNeumann1932}. HMT--MD
constructs a layer prior to that representation: a compatible section,
its observable route, transport memory, and the reader that publishes an
outcome. Hilbert space is not replaced; it appears as the codomain of
a declared linear realization.

The typed sequence is
\[
 \boxed{\begin{aligned}
 \text{genealogical section}
 &\longrightarrow \text{linear realization}
 \longrightarrow \text{phase and superposition}\\
 &\longrightarrow \text{observable and evolution}
 \longrightarrow \text{probability and measurement}
 \end{aligned}}
\]
Each arrow preserves its domain. In particular, the state vector does not
replace the TPK section, just as the real value does not replace
the numerical genealogy that produces it.

\subsection{Genealogical state and Hilbert realization}

\begin{definition}[Realized genealogical state]
A genealogical state is a quintuple
\begin{equation}
 \mathfrak G=(x_\infty,\Gamma,\mathcal F,\mathcal M,\mathcal R),
 \label{eq:estado-genealogico-realizado-rev2}
\end{equation}
where \(x_\infty\in\Omega_\infty\) is a compatible section,
\(\Gamma\) its observable route, \(\mathcal F\) the family of realizations,
\(\mathcal M\) the memory of the lifts, and \(\mathcal R\) the declared
reader. Let \(\mathsf P_{\rm TPK}\) be the category whose objects are the
prepared TPK sections, finite or coinductive, and whose morphisms are the
recorded routes compatible with their boundary data. In particular, the
depth restriction is a morphism
\[
 r_{n+1,n}:x_{n+1}\longrightarrow x_n.
\]
Let \(\mathsf{Hilb}_{\mathbb R,J}\) be the category whose objects
are pairs \((\mathcal H_{\mathbb R},J)\), with \(\mathcal H_{\mathbb R}\)
a real Hilbert space and \(J\) orthogonal, \(J^2=-I\), and whose morphisms are
bounded real-linear maps \(T\) intertwining the complex
structures, \(TJ_1=J_2T\). A quantum realization is a functor
\[
 \mathcal Q:\mathsf P_{\rm TPK}\longrightarrow
 \mathsf{Hilb}_{\mathbb R,J},
\]
and a \emph{pointed realization} is the pair \((\mathcal Q,\Psi)\) assigning
to each prepared section \(x\) a nonzero vector
\(\Psi_x\in\mathcal Q(x)\), compatible with the restrictions of the inverse
system.
\end{definition}

With the convention of linearity in the second argument, the associated Hermitian form
is written
\[
 \langle u,v\rangle_{\mathbb C}
 =\langle u,v\rangle_{\mathbb R}
  -i\langle u,J_E v\rangle_{\mathbb R}.
\]
Choosing the chart \(J_E\leftrightarrow i\) realizes the complex structure that
in Parts I--III is constructed as an internal root of \(-I\). For prefixes of
depth \(n\), set
\(\mathcal H_n:=\mathcal Q(x_n)\) and
\[
 p_{n+1,n}:=\mathcal Q(r_{n+1,n}):
 \mathcal H_{n+1}\longrightarrow\mathcal H_n.
\]
Compatibility between restriction and realization requires
\begin{equation}
 p_{n+1,n}\Psi_{n+1}(x_{n+1})=\Psi_n(x_n).
 \label{eq:compatibilidad-realizacion-hilbert-rev2}
\end{equation}
If \(r_{\infty,n}:x_\infty\to x_n\) is the coinductive restriction and
\(p_n:=\mathcal Q(r_{\infty,n}):\mathcal H_\infty\to\mathcal H_n\), a
limit \(\mathcal H_\infty\) realizes the complete section when
\(p_n\Psi_\infty(x_\infty)=\Psi_n(x_n)\) at every depth.

\begin{principle}[Realized quantum state]
The pure state associated with a section is represented by the ray of the
nonzero vector \(\Psi=\Psi_\infty(x_\infty)\); a statistical preparation is
represented by a positive trace-one density operator on the same
realization. Global-phase identification acts in the representation:
the section and its record remain in the genealogical domain.
\end{principle}

\subsection{Superposition and internal phase}

If \(\Psi,\Phi\in\mathcal H_\infty\) and
\(a=\alpha+i\beta\in\mathbb C\), the action of \(a\) in the real representation
is \(\alpha I+\beta J_E\). The combination \(a\Psi+b\Phi\) is formed only after
transporting the two alternatives to a common fiber. The genealogical lifts
of \cref{chap:genealogia-operaciones-rev2} take values in
\(\text{output}\times\text{memory}\); they belong to the TPK domain and are not
multiplied as operators. If a projective realization assigns unitary
operators \(U_g\) to visible transformations, its composition defect is
described separately by
\[
 U_gU_h=c(g,h)U_{gh},
 \qquad c(g,h)\in U(1),
\]
with \(c(e,g)=c(g,e)=1\) and
\[
 c(g,h)c(gh,k)=c(h,k)c(g,hk).
\]
A cocycle valued in another memory group acts on the realization only
after a representation of that group has been declared that intertwines the
operators \(U_g\).

\begin{principle}[Genealogical superposition]
Compatible alternatives are first transported in
\(\mathsf P_{\rm TPK}\) and are then combined linearly in
\(\mathcal H_\infty\). The action of \(J_E\) represents the relative phase, and
the route weight preserves the declared lifting datum.
\end{principle}

The interference is expressed through the cross terms.
\[
 \|\Psi_1+\Psi_2\|^2
 =\|\Psi_1\|^2+\|\Psi_2\|^2
  +2\operatorname{Re}\langle\Psi_1,\Psi_2\rangle.
\]
Two routes with identical endpoints may therefore contribute distinct phases
without their joint publication erasing the difference between their histories.

\subsection{Observables and Spectral Calculus}

Let \(O:\Omega_\infty\to Y\subseteq\mathbb R\) be a real genealogical reader. A quantum realization of
\(O\) is a self-adjoint operator \(\widehat O\), with dense domain
\(D(\widehat O)\), whose PVM satisfies the compatibility condition
\[
 P_O(B)\Psi_x
 \quad\hbox{represents the restriction of \(x\) to \(O^{-1}(B)\)}
\]
for each Borel set \(B\subseteq\mathbb R\) in the support of the preparation.
In finite dimension,
\[
 \widehat O=\sum_a o_a P_a,
 \qquad
 P_a P_b=\delta_{ab}P_a,
 \qquad
 \sum_a P_a=I;
\]
in general dimension, the sum is replaced by the spectral integral
\(\widehat O=\int\lambda\,\mathrm dP_O(\lambda)\).

\begin{principle}[Typed observable]
The physical observable is the pair formed by the reader in the
genealogical domain and its self-adjoint representation. Spectral calculus publishes the
possible results of the representation without erasing the operation that
defined them.
\end{principle}

For a normalized pure state \(\Psi\in D(\widehat O^2)\),
\[
 \langle O\rangle_\Psi
 =\langle\Psi,\widehat O\Psi\rangle,
 \qquad
 (\Delta_\Psi O)^2
 =\langle\Psi,(\widehat O-\langle O\rangle_\Psi)^2\Psi\rangle.
\]
The formulation by means of \(\operatorname{tr}(\rho\widehat O)\) extends the definition to mixed
states when \(\rho\widehat O\) is trace class; the second moment likewise requires
\(\rho\widehat O^2\) to be trace class.

\subsection{Evolution, Hamiltonian, and Schrödinger equation}

Genealogical evolution preserves composition of transports. Its
realization satisfies
\[
 U(t_3,t_2)U(t_2,t_1)=U(t_3,t_1),
 \qquad U(t,t)=I.
\]
When it preserves the inner product and satisfies
\(U(t)J_E=J_EU(t)\), it is orthogonal in the real chart and unitary in the
complex chart. For a strongly continuous one-parameter unitary group,
Stone's theorem provides a self-adjoint and \(J_E\)-linear
generator \(H\)~\cite{Stone1932}:
\[
 U(t)=\exp\!\left(-J_E H t/\hbar_{\rm int}\right).
\]
On the domain \(D(H)\), differentiation with respect to time gives
\begin{equation}
 \hbar_{\rm int}J_E\dot\Psi(t)=H\Psi(t),
 \label{eq:schrodinger-genealogico-rev2}
\end{equation}
and the chart \(J_E\leftrightarrow i\) produces
\(i\hbar_{\rm int}\dot\Psi=H\Psi\).

The calendar \(C_{108}\) realizes this correspondence in finite form:
the self-adjoint Hamiltonian constructed for it has, over \(t_0\),
the unitary one-iteration shift as its exponential. The
continuous group additionally requires strong continuity; it does not follow
from the periodicity of the calendar alone.

\subsection{Probability and quadratic representation of the branches}

Let \(\mathcal C_n=\{C_s:s\in\mathcal S_n\}\) be the partition of
\(\Omega_\infty\) by a finite set \(\mathcal S_n\) of prefixes of
depth \(n\), and let \(\mu\) be a compatible cylindrical probability
measure, \(\mu(\Omega_\infty)=1\). Let \(E\) be a real Hilbert space
with complex structure \(J_E\). On
\(\mathcal H_n=\ell^2(\mathcal S_n)\otimes E\), define
\begin{equation}
 \Psi_n=\sum_{s\in\mathcal S_n}
 \sqrt{\mu(C_s)}\,|s\rangle\otimes e^{\theta_s J_E}u_0,
 \label{eq:estado-amplitudes-medida-rev2}
\end{equation}
with \(u_0\) unitary.

\begin{theorem}[Quadratic representation of a branch measure]
\label{thm:born-cilindrico-rev2}
For a cylindrical event
\(A=\bigcup_{s\in I_A}C_s\), let
\[
 P_A=\left(\sum_{s\in I_A}|s\rangle\langle s|\right)\otimes I_E.
\]
Then
\[
 \|P_A\Psi_n\|^2=\mu(A).
\]
\end{theorem}

\begin{proof}
The prefix vectors are orthogonal and
\(e^{\theta_s J_E}\) preserves the norm. Therefore,
\[
 \|P_A\Psi_n\|^2
 =\sum_{s\in I_A}\mu(C_s)=\mu(A).
\]
\end{proof}

The quadratic rule is obtained here for the PVM of finite-depth cylindrical
events as an exact representation of the declared branch
measure~\cite{Born1926}. Conversely, Gleason's theorem
characterizes, in dimension at least three, noncontextual additive measures
on projectors through density operators~\cite{Gleason1957}. Coupling
with the apparatus and conditioning on the outcome are developed
subsequently in measurement theory.

\subsection{Incompatibility of readers and uncertainty}

Translations and nonadic phases define the Weyl pair.
\[
 \mathcal H_9=\ell^2(\mathbb Z/9\mathbb Z),
 \qquad
 \omega_9=e^{2\pi J_E/9},
 \qquad
 X|n\rangle=|n+1\rangle,
 \qquad
 Z|n\rangle=\omega_9^n|n\rangle.
\]
The indices are read modulo \(9\), and \(\omega_9\) is the complex scalar
realized by the structure \(J_E\) on \(\mathcal H_9\).
The identity
\[
 ZX=\omega_9XZ
\]
realizes the finite Heisenberg group of order \(9^3=729\). This cardinality
belongs to the group of phases and translations and retains its type against the
other occurrences of \(729\).

For self-adjoint observables \(A\) and \(B\), a normalized vector
\(\Psi\in D(AB)\cap D(BA)\) with finite second moments permits the application of
Cauchy--Schwarz to obtain the
Robertson--Heisenberg inequality
\begin{equation}
 \Delta_\Psi A\,\Delta_\Psi B
 \geq\frac12
 \bigl|\langle\Psi,[A,B]\Psi\rangle\bigr|.
 \label{eq:robertson-hmt-rev2}
\end{equation}
The incompatibility thus appears at coordinated levels: the commutator of
the APP evaluations, the Weyl relation in the representation, and the dispersion
bound for the physical readings~\cite{Heisenberg1927,Robertson1929}.

For \(\Psi\neq0\), the discrete Fourier transform supplies the support bound
\[
 |\operatorname{supp}\Psi|\,
 |\operatorname{supp}\widehat\Psi|\geq9,
\]
and, for mutually unbiased position and phase bases, the Shannon entropic
inequality, with natural logarithms,
\[
 H_Q(\Psi)+H_P(\Psi)\geq\log9
\]
offers a third exact formulation~\cite{MaassenUffink1988}. Here
\(H_Q\) and \(H_P\) are the Shannon entropies of the probability
distributions in the two bases.

\subsection{Tensor product of paths and entanglement generation via shared memory}

Genealogical composition glues sections over boundary data. Fix
on \(\mathsf P_{\rm TPK}\) a monoidal structure
\((\boxtimes,\mathbf 1,a,l,r)\), with declared product, unit, and associativity
and unit coherences. A strong monoidal realization
provides natural isomorphisms
\[
 \mu_{A,B}:
 \mathcal Q(A)\widehat\otimes\mathcal Q(B)
 \xrightarrow{\sim}\mathcal Q(A\boxtimes B).
\]
In the Hilbert image there appears
\(\mathcal H_A\otimes\mathcal H_B\). A nonseparable pure state admits a
Schmidt decomposition
\[
 |\Psi\rangle
 =\sum_r\sqrt{p_r}\,|r\rangle_A\otimes|r\rangle_B,
\]
and satisfies \(p_r\geq0\), \(\sum_rp_r=1\), with
\(\{|r\rangle_A\}\) and \(\{|r\rangle_B\}\) orthonormal Schmidt bases,
and its reductions have the same nonzero spectrum. The entropy associated with
that spectrum receives, in the Barbero--Immirzi chapter, an area and boundary
realization.

\begin{principle}[Genealogical composition]
The tensor product represents a gluing of sections over a shared
boundary. Entanglement expresses the nonfactorization of the
realized section with respect to the chosen bipartition.
\end{principle}

\subsection{Kinematic Reconstruction Theorem}

\begin{theorem}[Hilbertian Conditional Reconstruction of the Genealogical State]
\label{thm:cinematica-cuantica-tpk-rev2}
Let \(\Omega_\infty\) be a space of compatible sections, let
\((\mathsf P_{\rm TPK},\boxtimes,\mathbf 1)\) be the monoidal category of routes with the declared
coherences, and let
\[
 \mathcal Q:\mathsf P_{\rm TPK}\longrightarrow
 \mathsf{Hilb}_{\mathbb R,J}
\]
a strong monoidal realization, pointed by a compatible section
\(\Psi_x\in\mathcal Q(x)\). Suppose that:
\begin{enumerate}
 \item \(J_E^2=-I\), \(J_E\) preserves the inner product, and
the morphisms of \(\mathcal Q\) intertwine the complex structures;
 \item each declared real reader is represented by a self-adjoint operator
 whose PVM satisfies the compatibility condition with the reader's fibers;
 \item each continuous evolution is represented by a strongly continuous
 one-parameter unitary group \(U(t)\), with
 \(U(t)J_E=J_EU(t)\), and a self-adjoint \(J_E\)-linear generator \(H\);
 \item the preparation induces a compatible cylindrical probability measure
 and the state \(\Psi_n\) of \eqref{eq:estado-amplitudes-medida-rev2}.
\end{enumerate}
The complex chart of \(\mathcal Q\) then realizes states as rays or densities,
complex superposition, self-adjoint observables, unitary evolution generated
by \(H\), tensor composition, and quadratic probabilities for the
finite-depth cylindrical PVM.
\end{theorem}

\begin{proof}
The first hypothesis converts every real fiber into a complex Hilbert
space and makes the transports complex-linear. The second permits
application of the spectral theorem with the typed reader. The third supplies the
Hamiltonian through Stone's theorem and the differential equation on
\(D(H)\). The fourth and \cref{thm:born-cilindrico-rev2} give the quadratic
rule on cylindrical events. Strong monoidality realizes the
tensor product and its notion of separability.
\end{proof}

\begin{statusbox}[title={Result and scope}]
The Hilbertian reconstruction is coordinated as a property of a
realization of the genealogical state: the complex structure derives from
\(J_E^2=-I\); evolution, from the representation of routes; probability,
from a compatible measure; and entanglement, from composition over a
shared boundary. The theorem is exact under its four hypotheses and does not
automatically identify TPK with a concrete quantum theory. The
construction of a concrete physical realization preserves its sources:
Hamiltonian, apparatus, readers, domain, and boundary conditions.
\end{statusbox}
